
\begin{proposition}[The unrestricted irreducible zero-Euler criterion]\label{cdp:st:C.2.UN.12}
Let $f:Y\to B$ be a proper holomorphic map from a smooth
complex threefold to a smooth complex curve, with connected
fibres and smooth general fibre a complex two-torus. Fix an
irreducible full fibre $mS$, and write
$h=\eta\pi:Z\to T\to S$ for normalization followed by its
minimal resolution. Every such support satisfies $e(S)\ge0$.
Its Euler number is zero precisely in the following geometric
cases, stated for these actual maps.

A normal support is a smooth torus or a smooth bielliptic
surface, with multiplicities as in \ref{cdp:st:C.2.UN.5}.

A nonnormal support with irrational ruled normalization has
smooth minimal normalization ruled over an elliptic curve.
Its conductor image and reduced inverse image are disjoint
unions of smooth elliptic curves. The transversal types and
their finite conductor maps are exactly those in the table:
\[
\begin{array}{c|c|c}
 \text{normalization conductor}&\text{transversal type}
       &\text{degree over the image}\\ \hline
 H\text{ reduced, ruling degree }2&xy=0&2\\
 2C,\ C\text{ a section}&x^2-y^3=0&1\\
 2C,\ C\text{ a section}&x^2-y^4=0&2\\
 2C,\ C\text{ a section}&xy(x-y)=0&3
\end{array}                                                   \tag{C.IZ.1}\label{cdp:eq:C.IZ.1}
\]
Every displayed map of elliptic curves is unramified; the
degree-one map is an isomorphism. The transversal type is
constant along each conductor component. In the first row,
degree two means the total degree over each image component;
its inverse image may have two connected components. There are no
contracted horizontal sections in the zero-Euler case.

For a nonnormal support whose normalization resolution
has a Hopf surface as its minimal model, put $R=D_{\rm red}$,
where $D$ is the effective
resolution conductor with $K_Z+D=h^*\omega_S$. Let $r$ be
the number of rational components of $R$. Put
$C=\Gamma_{\rm red}$ for the conductor image, let $n_0$ be
the number of its rational-normalized components, and set
$v(C)=\sum_x(s_x-1)$, using the actual number $s_x$ of
normalization branches at each singular point. The exact
zero criterion is
\[
 b_2(Z)=r,\qquad
 \begin{cases}
 n_0=0,\quad v(C)=0,\quad\sum_x\delta_x\le1,\quad\text{or}\\
 n_0=1,\quad v(C)=2,\quad\sum_x\delta_x=2.
 \end{cases}                                                   \tag{C.IZ.2}\label{cdp:eq:C.IZ.2}
\]
Every other normalized conductor component is elliptic.
In the second line every singular point satisfies
$\delta_x=s_x-1$. No normalization singularity lies outside
the conductor in either line. An irreducible support with
rational normalization resolution always has $e(S)\ge2$.

This is an equivalence for fibres of the specified map.
The displayed curve data retain their normalization and
resolution morphisms; no assertion of smoothability of
arbitrary abstract surface gluings is included.
\end{proposition}
\begin{proof}
We first check the classification's scope without an ambient
Chern-class assumption. On a sufficiently small fibre tube
$U=f^{-1}(\Delta)$, the evaluation of a generator of the
rank-one torsion-free sheaf $f_*K_U$ gives
$K_U\simeq\OO_U(kS)$, as proved with its divisor comparison
in \ref{cdp:st:C.2.UN.7}. Since $mS$ is principal, $K_U$ and
$\omega_S=\OO_U((k+1)S)|_S$ are torsion. The actual
$\OO_{mS}$ filtration and flatness give
$\chi(\OO_S)=0$. For nonnormal $S$, finite duality and
resolution adjunction give the nonempty effective divisor
$D$ and the stated torsion line $K_Z+D$.

Every positive pluricanonical section on $Z$ would, after
taking a power trivializing that torsion line, give a
nonzero section of $\OO_Z(-aD)$ for some $a>0$.
Its inclusion in $\OO_Z$ would be a nonzero holomorphic
function on the compact connected surface vanishing along
$D$, a contradiction. Thus $\kappa(Z)=-\infty$.
The surface classification, together with the actual
minimal-model blowdown, makes $Z$ rational, irrational
ruled, or class VII. In the last case \ref{cdp:st:C.2.UN.19} proves,
through the specified cyclic cover and quotient maps,
that its minimal model is Hopf. Thus the formulation
above includes every class VII possibility for the
actual degeneration. Normal supports have already been
classified with their finite root covers in \ref{cdp:st:C.2.UN.5}.

\ref{cdp:st:C.2.UN.7} proves strict positivity for rational $Z$.
It also proves strict positivity when an irrational
horizontal section is contracted, and extends every
noncontracted irrational ruled comparison to this tube.
The precise replacements for the earlier ambient $b_2=0$
uses are the torsion canonical and normal lines, the
vanishing of $\chi(\OO_S)$, and
$\int_Zh^*c_2(T_U)=0$, all proved there through restriction
to the nearby torus.

For clarity, the horizontal possibilities exhaust the
ruling degree two of $D$. A reduced horizontal part is
covered by C.2.IRR.1: its adjunction comparison makes its
curves smooth elliptic and disjoint, its normalization
smooth minimal ruled, and its conductor map an unramified
double cover with product-node germs. A nonreduced
horizontal part is $2C$ for a section $C$. At a generic
conductor point, the plane-curve conductor exponent of a
normalization branch equals its intrinsic conductor plus
its intersections with the other branches. Exponent two
permits one cusp branch, two smooth branches meeting with
multiplicity two, or three pairwise transverse smooth
branches. These are the last three rows of \eqref{cdp:eq:C.IZ.1}.
An additional branch would give exponent at least three.
The exact cusp, tacnode and cubic-cover comparisons in
\ref{cdp:st:A.5.20}, \ref{cdp:st:A.5.18} and \ref{cdp:st:C.2.D4.1}--2, extended in \ref{cdp:st:C.2.UN.7},
give the asserted elliptic image, constant type and
degrees one, two and three. A vertical rational curve
cannot dominate that elliptic image, so every generic
branch is accounted for by the specified horizontal map.
All these comparisons remove vertical conductor and
normalization exceptional curves. Their Euler formula is
\[
 e(S)=e(T)-e((\eta^{-1}C)_{\rm red})+e(C)=0,             \tag{C.IZ.3}\label{cdp:eq:C.IZ.3}
\]
since each term vanishes. This proves both directions of
the ruled part of the criterion.

It remains to check and extend the class VII part. The
proof of \ref{cdp:st:C.2.VII.1} uses only the actual torsion relation
$K_Z+D$; it gives
\[
 e(R)=r,\qquad r\le b_2(Z).                              \tag{C.IZ.4}\label{cdp:eq:C.IZ.4}
\]
The proof of \ref{cdp:st:C.2.VII.2} uses only this relation,
$\omega_S$ torsion and $\chi(\OO_S)=0$. They have just
been established on $U$, so its finite-duality and
curve-normalization maps give
\[
 \chi(\OO_C)\ge-1,\qquad e(C)\ge0.                      \tag{C.IZ.5}\label{cdp:eq:C.IZ.5}
\]
Its normalizations are rational or elliptic, because they
are finite images of the rational or elliptic curves in
$R$. These statements have no remaining ambient Betti
assumption.

We sharpen the exceptional terms in the earlier comparison.
At a point of $T$ outside the conductor, $\eta$ is an
isomorphism. A singular point there would consequently be
an isolated normal hypersurface singularity of $S$ in
the smooth tube. The isolated Milnor-fibre argument
\ref{cdp:st:A.5.17}, retaining the local unit and multiplicity of
the full fibre, excludes that point. Thus $T$ is smooth
outside the conductor. Every resolution-exceptional curve
of coefficient zero in $D$ is disjoint from $R$: minimal
resolution adjunction gives $K_Z\cdot E\ge0$, while
$D\cdot E\ge0$ and $(K_Z+D)\cdot E=0$ force both
numbers to vanish. Each exceptional fibre over the
conductor meets its strict transform and is connected.
Hence it cannot contain such a disjoint curve, and its
whole inverse image lies in $R$. There are no exceptional
fibres outside the conductor, by the preceding local
argument. We have therefore identified the actual
complement isomorphism
\[
 Z\setminus R\xrightarrow{\ \simeq\ }
 T\setminus(\eta^{-1}C)_{\rm red}
 \xrightarrow{\ \simeq\ }S\setminus C.                \tag{C.IZ.6}\label{cdp:eq:C.IZ.6}
\]
Its first map is the resolution restricted to this open
set, and its second is the finite normalization restricted
where its conductor ideal is the unit ideal.

A class VII surface and its point blowups have
$b_1=b_3=1$, so $e(Z)=b_2(Z)$. Euler additivity in
\eqref{cdp:eq:C.IZ.6} now yields the exact formula
\[
 e(S)=b_2(Z)-r+e(C),\qquad
 b_2(Z)-r\ge0,\quad e(C)\ge0.                           \tag{C.IZ.7}\label{cdp:eq:C.IZ.7}
\]
The curve normalization sequence identifies the remaining
term, rather than replacing its singular points by a
model:
\[
 e(C)=2n_0-v(C),\qquad
 \chi(\OO_C)=n_0-\sum_x\delta_x,\qquad
 \delta_x\ge s_x-1.                                    \tag{C.IZ.8}\label{cdp:eq:C.IZ.8}
\]
Equations \eqref{cdp:eq:C.IZ.5} and \eqref{cdp:eq:C.IZ.8} imply $e(C)\ge n_0-1$.
Consequently $e(C)=0$ allows exactly $n_0=0,v=0$ or
$n_0=1,v=2$. In the first case the coherent bound gives
$\sum\delta_x\le1$. In the second, that bound and
$\delta_x\ge s_x-1$ give $\sum\delta_x=2$, with equality
at every singular point. Conversely either of these
two curve patterns makes $e(C)=0$ in the exact
topological formula. Equation \eqref{cdp:eq:C.IZ.7} proves precisely
\eqref{cdp:eq:C.IZ.2}, and proves the sign for every class VII
support. All three nonnormal surface classes have now
been included, which proves the assertion.
\end{proof}

The class VII surface input is the proved anticanonical
configuration theorem checked in \ref{cdp:st:C.2.VII.1}, with its exact
references to \cite{Dloussky2006,DOT2003}. \ref{cdp:st:C.2.UN.19} further
identifies the minimal model in this branch as Hopf;
the existing conductor comparison applies through its
actual point-blowup map. The result here
transfers that input to arbitrary irreducible torus fibres
through the canonical evaluation, isolated-singularity
exclusion and the complement map \eqref{cdp:eq:C.IZ.6}. The unreduced
fibre multiplicity is retained throughout. The remaining
unrestricted classification problem in C.2 concerns
reducible supports.


